\section{LOB Recap (Prior Work)}
\label{sec:lob-recap}

In our prior work \cite{solovev2026lob}
(figshare DOI \href{https://doi.org/10.6084/m9.figshare.31859557}%
{10.6084/m9.figshare.31859557}) we proposed DA-BiGRU-CNN for LOB
mid-price prediction.  Two BiGRU branches operate on price-series and
volume-series features respectively, with a multi-scale Conv1d fusion
head; a 21-dimensional microstructure-feature embedding modulates the
fusion stage.  The composite loss of \Cref{eq:composite-loss} directly
optimises weighted-Pearson correlation, which is also the evaluation
metric reported on the held-out test fold.

Three findings from the LOB study motivate the cross-domain transfer we
study here.  First, the dual-branch GRU on a compact set of 53
hand-selected microstructure features achieves
weighted-Pearson $\LobGRUwPearson{}$ on test, outperforming a
LightGBM \cite{ke2017lightgbm} baseline trained on 205 features
($\LobLGBMwPearson{}$) by $\LobGRUOverLGBM{}$.  Second, feature
sufficiency is asymmetric across model families: expanding 53 features
to 219 improves the GRU by only $\LobFeatureSuffGRU$ while the LightGBM
gains $\LobFeatureSuffLGBM$ over the same extension --- the
domain-aware branches absorb the structural information at low feature
dimensionality.  Third, a naive GRU-plus-gradient-boosting ensemble
\emph{degrades} performance by $0.05$--$1.6\%$ against the GRU alone:
the two model families fail to be complementary on this task, again
suggesting the GRU has already captured the available signal.
\Cref{tab:lob-recap} reproduces the headline numbers.

\begin{table}[t]
\centering
\caption{LOB results from the prior work \cite{solovev2026lob}.
Test weighted-Pearson on the held-out fold; $\Delta$ is the relative
improvement over the GRU.}
\label{tab:lob-recap}
\begin{tabular}{lccc}
\toprule
Model & \#Features & wPearson & $\Delta$ vs GRU \\
\midrule
DA-BiGRU-CNN (ours) & 53  & $\LobGRUwPearson$  & --- \\
LightGBM            & 205 & $\LobLGBMwPearson$ & $-\LobGRUOverLGBM$ \\
GRU + LightGBM ensemble & 53 + 205 & 0.262 & $-1.6\%$ \\
\bottomrule
\end{tabular}
\end{table}

These findings constitute the baseline methodology we now test in a
structurally different domain (\Cref{sec:methodology-defi,%
sec:defi-experiment}): the same architecture, the same composite loss,
and a re-identified domain-natural decomposition pair --- substituting
$(\varepsilon, u+\text{ctx})$ for $(\text{price}, \text{volume})$.
